% =========================
% report/implementation.tex  (REPLACEMENT FOR VAE + cWGAN-GP REPORT)
% NOTE: main.tex already contains \section{Implementation}; do NOT add another \section here.
% =========================
\label{sec:implementation}

\subsection{Objective and Scope}
This section documents a complete, reproducible implementation of:
(i) baseline and referenced VAEs (including enhanced-capacity objective and hierarchical latent architecture) trained on Dynamic MNIST, and
(ii) a conditional WGAN-GP with attention and spectral normalization trained on FashionMNIST for data augmentation, followed by classifier-based evaluation.

All experiments produce audit-ready artifacts: training curves (loss components), generated samples at fixed checkpoints, latent diagnostics (posterior statistics, activity measures, traversals/PCA), and classifier metrics under real-only and augmented regimes.

\subsection{Repository Layout and Script Responsibilities}
A clean separation of concerns is used to ensure traceability:
\begin{verbatim}
project_root/
  src/
    vae/
      01_train_baseline_vae.py
      02_train_proposed_vae.py
      03_eval_latent_diagnostics.py
      04_estimate_likelihood.py
      05_visualize_pseudo_inputs.py
    gan/
      01_build_splits.py
      02_train_cwgan_gp.py
      03_generate_augmented_set.py
      04_train_classifier.py
      05_compare_results.py
  data/
    mnist_dynamic/      (or generated on-the-fly from MNIST with dynamic binarization)
    fashionmnist/
    splits/             (indices for non-overlapping subsets)
  results/
    vae/
      curves/
      samples/
      latents/
      pseudo_inputs/
      likelihood/
    gan/
      curves/
      samples/
      augmented_data/
      classifier/
  report/
    notation.tex
    implementation.tex
    experiments.tex
    evaluation.tex
    reproducibility.tex
    references.bib
\end{verbatim}

\subsection{Software Stack and Determinism Controls}
\paragraph{Core dependencies.}
\begin{itemize}
  \item PyTorch for model definition, training, and automatic differentiation.
  \item torchvision for MNIST/FashionMNIST loading and standard transforms.
  \item NumPy for numerics and dataset index management.
  \item Matplotlib for report-quality plots.
\end{itemize}

\paragraph{Reproducibility controls.}
All scripts:
\begin{itemize}
  \item fix random seeds (Python/NumPy/PyTorch),
  \item log configurations (YAML/JSON) alongside outputs,
  \item record device (CPU/GPU), PyTorch version, and commit hash (if applicable),
  \item store dataset split indices to guarantee non-overlapping subsets.
\end{itemize}

\subsection{Part I: VAE Implementations (Dynamic MNIST)}
\subsubsection{Dynamic MNIST data pipeline}
Dynamic MNIST is implemented via \emph{dynamic binarization}: at each epoch/iteration, grayscale pixel intensities (in $[0,1]$) are treated as Bernoulli parameters and binarized by sampling. This prevents overfitting to a fixed binarization and is standard for likelihood-based comparisons.

\subsubsection{Baseline VAE}
\paragraph{Model parameterization.}
We use an encoder $q_\phi(z\mid x)=\mathcal{N}(\mu_\phi(x),\mathrm{diag}(\sigma_\phi^2(x)))$ and a Bernoulli decoder $p_\theta(x\mid z)$ parameterized by a neural network producing logits (or probabilities after sigmoid).

\paragraph{Objective and logging.}
The baseline optimizes:
\[
\mathcal{L}_{\mathrm{ELBO}}(x)=
\mathbb{E}_{q_\phi(z\mid x)}[\log p_\theta(x\mid z)]
-
\mathrm{KL}(q_\phi(z\mid x)\|p(z)).
\]
We log $\mathcal{L}_{\mathrm{rec}}$ and $\mathcal{L}_{\mathrm{KL}}$ separately per step/epoch, along with total loss.

\subsubsection{Proposed VAE: enhanced-capacity objective and hierarchical structure}
\paragraph{Enhanced-capacity objective (learned mixture prior).}
We implement the paper’s modification by replacing $p(z)$ with a flexible prior $p_\lambda(z)$ induced by pseudo-inputs $\{u_k\}$:
\[
p_\lambda(z)=\frac{1}{K}\sum_{k=1}^K q_\phi(z\mid u_k),
\qquad
\mathcal{L}(x)=
\mathbb{E}_{q_\phi(z\mid x)}[\log p_\theta(x\mid z)]
-
\mathrm{KL}(q_\phi(z\mid x)\|p_\lambda(z)).
\]
Pseudo-inputs are learned jointly with $\theta,\phi$ and visualized during training.

\paragraph{Hierarchical latent implementation.}
We implement a two-level hierarchy ($z_2\rightarrow z_1\rightarrow x$) as specified in the assignment/paper:
\[
p_\theta(x,z_1,z_2)=p_\theta(x\mid z_1)\,p_\theta(z_1\mid z_2)\,p(z_2),
\quad
q_\phi(z_1,z_2\mid x)=q_\phi(z_2\mid x)\,q_\phi(z_1\mid x,z_2).
\]
We report KL terms per level to analyze latent utilization and posterior collapse behavior.

\subsubsection{Latent diagnostics and posterior collapse analysis}
We compute:
\begin{itemize}
  \item posterior mean/variance summaries across the dataset,
  \item per-dimension KL contributions and the number of \emph{active units},
  \item PCA on latent codes and controlled latent traversals,
  \item pseudo-input visualizations (inputs and their induced latent modes),
  \item likelihood estimates using importance sampling (IWAE-style).
\end{itemize}

\subsection{Part II: Conditional WGAN-GP for Augmentation (FashionMNIST)}
\subsubsection{Non-overlapping data splits}
We construct two disjoint subsets:
(i) a GAN subset for learning $p(x\mid y)$, and
(ii) a classifier subset for downstream evaluation.
Per-class sampling is implemented exactly as required (e.g., $N$ per class) and split indices are saved under \texttt{data/splits/}.

\subsubsection{Architectural modules: SN, Self-Attention, Channel Attention}
\paragraph{Spectral Normalization (SN).}
We wrap Conv2d/ConvTranspose2d layers with SN to control the spectral norm of weights, improving critic stability.

\paragraph{Self-Attention (SA).}
SA modules follow the assignment’s tensor naming and shapes ($f,g,h,v$), compute attention maps of size $(B,N,N)$ with $N=HW$, and apply a residual connection $y=\gamma o + x$ with learnable scalar $\gamma$.

\paragraph{Channel Attention (CA).}
CA uses global average pooling, an MLP (two FC layers with ReLU), and sigmoid gating to recalibrate feature channels.

\subsubsection{Conditional generator and projection discriminator}
\paragraph{Conditional generator.}
Class embeddings are concatenated with noise: $\tilde{z}=[z; e_G(y)]$, then mapped through the specified upsampling blocks to generate $x_{\mathrm{fake}}$.

\paragraph{Projection discriminator.}
The critic computes
\[
D_\omega(x,y)=s_\omega(\phi_\omega(x))+\langle \phi_\omega(x), e_D(y)\rangle,
\]
ensuring label conditioning via inner products in feature space.

\subsubsection{WGAN-GP losses and training loop}
\paragraph{Loss terms.}
We implement the assignment’s exact WGAN-GP loss with:
\begin{itemize}
  \item Wasserstein term,
  \item gradient penalty ($\lambda_{\mathrm{GP}}=10$),
  \item drift penalty ($0.001$),
  \item embedding $L_2$ penalty ($0.001$).
\end{itemize}

\paragraph{Optimization protocol.}
We follow the specified hyperparameters:
Adam with $\mathrm{LR}(G)=2\times10^{-4}$, $\mathrm{LR}(D)=4\times10^{-4}$, betas $(0.0,0.9)$; scheduler $\gamma=0.95$ every 100 steps; $N_{\mathrm{critic}}=5$ critic steps per generator step; fixed noise/labels for checkpoint grids at steps $\{0,1000,2000,3000,4000\}$.

\subsubsection{Augmentation and classifier evaluation}
We generate a controlled number of synthetic samples per class from the trained generator, build an augmented training set (real + synthetic), and train the same classifier architecture under both regimes. Evaluation is performed on the unchanged FashionMNIST test set. We report overall accuracy and class-wise performance, supported by confusion matrices and per-class error rates where relevant.



% =========================
% report/experiments.tex  (REPLACEMENT FOR VAE + cWGAN-GP REPORT)
% NOTE: main.tex already has a Data/Protocol section; keep this file at \subsection level.
% =========================

\subsection{Datasets}
\paragraph{Dynamic MNIST.}
Dynamic MNIST is derived from MNIST by sampling a binarized image at each iteration using pixel intensities as Bernoulli parameters. This yields a stochastic training set that prevents models from overfitting a fixed binarization and is standard for likelihood-based evaluation.

\paragraph{FashionMNIST (limited-data regime).}
FashionMNIST is used for conditional generation and augmentation. All images are normalized using:
\[
x \leftarrow \frac{x-0.5}{0.5},
\]
matching the assignment’s transform specification.

\subsection{Train/Validation/Test Protocol}
\paragraph{VAE.}
VAEs are trained on the Dynamic MNIST training split. Where needed, a validation subset is used for early diagnostics (without altering test evaluation). Final reporting uses held-out test evaluation for likelihood estimation and latent analysis.

\paragraph{GAN augmentation.}
The FashionMNIST training set is split into two \emph{non-overlapping} subsets:
\begin{itemize}
  \item GAN subset: used exclusively to train the conditional WGAN-GP.
  \item Classifier subset: used to train the downstream classifier (baseline and augmented).
\end{itemize}
The split is stratified by class and uses the exact per-class sampling required by the assignment. Split indices are saved for reproducibility.

\subsection{Experimental Settings}
\paragraph{VAE experiments.}
We run:
\begin{itemize}
  \item baseline VAE with low-dimensional $d$,
  \item proposed model (enhanced objective + hierarchy) with the same $d$,
  \item proposed model with high-dimensional latent (e.g., $d=40$ as required).
\end{itemize}
For each run we export: training curves (reconstruction/KL/total), generated samples, posterior statistics, active units, PCA/traversals, pseudo-input visualizations, and likelihood estimates.

\paragraph{GAN experiments.}
We train the conditional WGAN-GP for $\mathrm{RUN}=4000$ steps with $N_{\mathrm{critic}}=5$ and log:
critic loss, generator loss, Wasserstein estimate, gradient penalty, drift term, and embedding regularization. Samples are saved at steps $\{0,1000,2000,3000,4000\}$ using fixed noise and fixed labels.

\subsection{Artifacts and Output Contract}
All experiments write:
\begin{itemize}
  \item plots under \texttt{results/*/curves/},
  \item samples/grids under \texttt{results/*/samples/},
  \item latent diagnostics under \texttt{results/vae/latents/},
  \item pseudo-input visualizations under \texttt{results/vae/pseudo\_inputs/},
  \item classifier metrics under \texttt{results/gan/classifier/}.
\end{itemize}

